\documentclass[11pt]{article}

\usepackage[a4paper,margin=1.1in]{geometry}
\usepackage[T1]{fontenc}
\usepackage[utf8]{inputenc}
\usepackage{lmodern}
\usepackage{microtype}
\usepackage{amsmath,amssymb,amsthm,mathtools}
\usepackage{booktabs}
\usepackage{enumitem}
\usepackage{xcolor}
\usepackage[colorlinks=true,linkcolor=blue!50!black,citecolor=blue!50!black,urlcolor=blue!50!black]{hyperref}

\numberwithin{equation}{section}

\newtheorem{theorem}{Theorem}[section]
\newtheorem{proposition}[theorem]{Proposition}
\newtheorem{lemma}[theorem]{Lemma}
\newtheorem{corollary}[theorem]{Corollary}
\newtheorem{definition}[theorem]{Definition}
\newtheorem{assumption}[theorem]{Assumption}
\newtheorem{remark}[theorem]{Remark}
\newtheorem{example}[theorem]{Example}

\newcommand{\Sym}{\mathfrak S}
\newcommand{\cyc}{\#}
\newcommand{\C}{\mathbb C}
\newcommand{\gammacyc}{\gamma_n}
\newcommand{\dplus}{\delta_+}
\newcommand{\dminus}{\delta_-}
\newcommand{\Dtot}{\Delta}
\newcommand{\rank}{\mathcal R}
\newcommand{\Count}{\mathcal C}
\newcommand{\Envelope}{\mathcal E}
\newcommand{\Qloss}{\mathcal Q}

\title{A Conditional Bidefect Reduction for Aubrun's Partial-Transpose Moment Count}
\author{Kieran McShane}
\date{\today}

\begin{document}

\maketitle

\begin{abstract}
We isolate a purely combinatorial reduction in Aubrun's moment method for
partially transposed Wishart matrices.  The reduction replaces a one-parameter
defect count by a two-parameter bidefect count attached to the two Cayley
geodesic inequalities for the long cycle.  This separation is exactly what is
needed for the two analytic bases in the partial-transpose expansion: the
plus-defect is paid by a shifted \(d\)-base and the minus-defect is paid by a
second shifted base, instantiated by the \(\sqrt{s}\)-scale in the Wishart
application.  We prove the scalar absorption and the
genus-induction reduction in full detail.  The final enumerative input is
reduced to two familiar ingredients: the genus-zero one-sided geodesic count
and a Chapuy-type slicing recurrence for the associated unicellular maps.
The main result is a self-contained conditional reduction theorem.
\end{abstract}

\tableofcontents

\section{Introduction}

Let \(W\) be a complex Wishart matrix on
\(\C^d\otimes \C^d\), and let \(W^\Gamma\) denote the partial transpose on one
tensor factor.  Aubrun's proof of the sharp PPT threshold studies moments of
\(W^\Gamma\) by expanding them into permutation-indexed Wick terms
\cite{Aubrun2012}.  In the balanced regime \(s\asymp d^2\), the leading
terms are controlled by Cayley-geodesic permutations.  The delicate part is
not the rank budget itself; it is the count of the almost-geodesic terms and
the way their defects are absorbed by the two different analytic bases.

The purpose of this note is to state and prove the clean reduction that
emerges from that bookkeeping.  All Cayley lengths below are taken in the
Cayley graph of \(\Sym_n\) generated by transpositions.  Write \(n=m+1\), let
\(\gammacyc=(0\,1\,\cdots\,n-1)\in\Sym_n\), and let \(\cyc(\sigma)\) be the
number of cycles of a permutation, fixed points included.  For
\(\pi\in\Sym_n\), define the two oriented Cayley defects
\[
  \dplus(\pi)=n+1-\bigl(\cyc(\pi)+\cyc(\gammacyc\pi)\bigr),
  \qquad
  \dminus(\pi)=n+1-\bigl(\cyc(\pi)+\cyc(\pi\gammacyc^{-1})\bigr).
\]
The central counting object is therefore not a total-defect slice, but the
bidefect slice
\[
  \Count_m(a,b)
  =
  \#\{\pi\in\Sym_{m+1}:\dplus(\pi)=a,\ \dminus(\pi)=b\}.
\]
The analytic argument requires a bidefect bound of the following shape, made
precise in Definition~\ref{def:bidefect-envelope} and
Assumption~\ref{ass:count}:
\[
  \Count_m(a,b)
  \le
  4^{m+1}
  \binom{m+3}{a}
  \binom{m+1}{b}
  \Qloss_m^{a+b},
  \qquad
  \Qloss_m=(4(m+1))^{36},
\]
on the analytic ranges \(0\le a\le m+3\), \(0\le b\le m+1\).
The binomial factors are not decorative.  They are exactly what lets the
double sum split into two binomial expansions:
\[
\sum_{a,b}\Count_m(a,b)d^{m+3-a}S^{m+1-b}
\lesssim
4^{m+1}(d+\Qloss_m)^{m+3}(S+\Qloss_m)^{m+1}.
\]
This is the scalar mechanism needed by the partial-transpose moment estimate.
The rest of the note reduces this bidefect bound to two enumerative inputs:
the genus-zero count and a Chapuy-type slicing recurrence.  The precise form
of the reduction is Theorem~\ref{thm:main}.

\section{Permutation Defects}

Throughout the paper \(n=m+1\ge1\), and
\(\gammacyc=(0\,1\,\ldots\,n-1)\in\Sym_n\).  We write
\(\cyc(\sigma)\) for the number of cycles of \(\sigma\), including fixed
points, and
\[
  |\sigma|=n-\cyc(\sigma)
\]
for the Cayley length with respect to transpositions.

\begin{definition}[Oriented defects]
For \(\pi\in\Sym_n\), define
\[
  \dplus(\pi)=n+1-\cyc(\pi)-\cyc(\gammacyc\pi),
  \qquad
  \dminus(\pi)=n+1-\cyc(\pi)-\cyc(\pi\gammacyc^{-1}).
\]
The total defect is
\[
  \Dtot(\pi)=\dplus(\pi)+\dminus(\pi).
\]
\end{definition}

\begin{lemma}[Non-negativity]
For every \(\pi\in\Sym_n\),
\[
  \dplus(\pi)\ge0,
  \qquad
  \dminus(\pi)\ge0.
\]
\end{lemma}

\begin{proof}
The Cayley length is the minimal number of transpositions needed to write a
permutation, so it is subadditive by concatenating factorizations.  Applying
this triangle inequality gives
\[
  |\pi|+|\pi^{-1}\gammacyc^{-1}|\ge |\gammacyc^{-1}|=n-1.
\]
Since \(|\sigma|=n-\cyc(\sigma)\) and
\(\cyc(\pi^{-1}\gammacyc^{-1})=\cyc(\gammacyc\pi)\), this is equivalent to
\[
  \cyc(\pi)+\cyc(\gammacyc\pi)\le n+1.
\]
This proves \(\dplus(\pi)\ge0\).  For the minus defect, the same triangle
inequality gives
\[
  |\pi|+|\pi^{-1}\gammacyc|\ge |\gammacyc|=n-1.
\]
Equivalently,
\[
  (n-\cyc(\pi))+(n-\cyc(\pi^{-1}\gammacyc))\ge n-1.
\]
Finally
\[
  \cyc(\pi^{-1}\gammacyc)
  =
  \cyc((\pi^{-1}\gammacyc)^{-1})
  =
  \cyc(\gammacyc^{-1}\pi)
  =
  \cyc(\pi\gammacyc^{-1}),
\]
where the last equality uses that \(uv\) and \(vu\) have the same cycle
count.  Hence
\[
  \cyc(\pi)+\cyc(\pi\gammacyc^{-1})\le n+1,
\]
which is exactly \(\dminus(\pi)\ge0\).
\end{proof}

\begin{definition}[Aubrun rank]
The permutation rank entering the balanced Wick term is
\[
  \rank(\pi)=2\cyc(\pi)+\cyc(\gammacyc\pi)+\cyc(\pi\gammacyc^{-1}).
\]
\end{definition}

\begin{proposition}[Rank and bidefect identity]\label{prop:rank}
For every \(\pi\in\Sym_n\),
\[
  \rank(\pi)+\dplus(\pi)+\dminus(\pi)=2n+2.
\]
In particular \(\rank(\pi)\le 2n+2\).
\end{proposition}

\begin{proof}
This is immediate from the definitions:
\[
\begin{aligned}
  \dplus(\pi)+\dminus(\pi)
  &=
  2n+2
  -2\cyc(\pi)
  -\cyc(\gammacyc\pi)
  -\cyc(\pi\gammacyc^{-1}) \\
  &=2n+2-\rank(\pi).
\end{aligned}
\]
The last assertion follows from the non-negativity of the two oriented
defects.
\end{proof}

\begin{example}[The case \(n=3\)]
Let \(\gammacyc=(0\,1\,2)\).  The following sample rows illustrate the
identity in Proposition~\ref{prop:rank}.
\[
\begin{array}{c|cccccc}
\pi & \cyc(\pi) & \cyc(\gammacyc\pi) & \cyc(\pi\gammacyc^{-1})
& \dplus(\pi) & \dminus(\pi) & \rank(\pi)\\
\hline
\mathrm{id} & 3 & 1 & 1 & 0 & 0 & 8\\
(0\,1) & 2 & 2 & 2 & 0 & 0 & 8\\
\gammacyc & 1 & 1 & 3 & 2 & 0 & 6\\
\gammacyc^{-1} & 1 & 3 & 1 & 0 & 2 & 6
\end{array}
\]
In every row \(\rank+\dplus+\dminus=8=2n+2\).
\end{example}

\section{The Bidefect Count and Scalar Absorption}

The following definition is the exact count required by the shifted-base
moment estimate.

\begin{definition}[Bidefect count]\label{def:bidefect-envelope}
For \(a,b\ge0\), set
\[
  \Count_m(a,b)=
  \#\{\pi\in\Sym_{m+1}:\dplus(\pi)=a,\ \dminus(\pi)=b\}.
\]
Let
\[
  \Qloss_m=(4(m+1))^{36}.
\]
The exponent \(36\) is much larger than the genus induction alone requires;
it is retained here because it is the uniform polynomial loss used in the
surrounding compatible-couple estimates.  This note does not derive that
ambient loss: the scalar argument below only needs a loss parameter large
enough to dominate the genus-slicing factor, and it allows any larger
replacement \(Q\ge \Qloss_m\).
Define the bidefect envelope
\[
  \Envelope_m(a,b)=
  4^{m+1}
  \binom{m+3}{a}
  \binom{m+1}{b}
  \Qloss_m^{a+b}.
\]
The binomial factors are inserted so that the scalar sum factors into the two
shifted bases \((d+Q)^{m+3}\) and \((S+Q)^{m+1}\).  Since the reduction below
uses only a one-sided plus-defect genus count, the factor
\(\binom{m+1}{b}\Qloss_m^b\) is deliberate padding rather than independent
minus-defect information.
\end{definition}

The cutoffs in the next assumption are dictated by
Proposition~\ref{prop:rank}.  If \(\dplus(\pi)=a\) and
\(\dminus(\pi)=b\), then
\[
  \rank(\pi)=2m+4-a-b.
\]
In the partial-transpose Wick expansion, the offsets \(m+3\) and \(m+1\) are
the base exponents attached to the two separated parameters before defects
are charged.  Thus the separated bases have schematic powers
\[
  d^{m+3-a}S^{m+1-b}.
\]
The shifted-base majorant only needs the terms whose displayed exponents are
nonnegative.  Thus the analytic ranges are precisely
\[
  0\le a\le m+3,\qquad 0\le b\le m+1.
\]
The theorem below is a reduction for this analytic part of the sum; it makes
no claim about slices outside these ranges.

\begin{assumption}[Bidefect count bound]\label{ass:count}
For every \(m\ge0\), every \(0\le a\le m+3\), and every
\(0\le b\le m+1\),
\[
  \Count_m(a,b)\le \Envelope_m(a,b).
\]
\end{assumption}

In the scalar theorem, \(S\) is only a placeholder for the second analytic
base.  In the balanced Wishart application one substitutes the corresponding
\(\sqrt{s}\)-scale, but the absorption argument itself is just the displayed
two-variable inequality.

\begin{theorem}[Shifted-base scalar absorption]\label{thm:scalar}
Assume the bidefect count bound for a fixed \(m\).  Let \(d,S,Q\ge0\) and
\(\Qloss_m\le Q\).  Then
\[
\sum_{a=0}^{m+3}\sum_{b=0}^{m+1}
  \Count_m(a,b)d^{m+3-a}S^{m+1-b}
  \le
  4^{m+1}(d+Q)^{m+3}(S+Q)^{m+1}.
\]
\end{theorem}

\begin{proof}
By Assumption~\ref{ass:count}, the left-hand side is at most
\[
  4^{m+1}
  \sum_{a=0}^{m+3}\sum_{b=0}^{m+1}
  \binom{m+3}{a}\binom{m+1}{b}
  \Qloss_m^{a+b}
  d^{m+3-a}S^{m+1-b}.
\]
The double sum factors:
\[
\begin{aligned}
&4^{m+1}
 \left(\sum_{a=0}^{m+3}
   \binom{m+3}{a}\Qloss_m^a d^{m+3-a}\right)
 \left(\sum_{b=0}^{m+1}
   \binom{m+1}{b}\Qloss_m^b S^{m+1-b}\right) \\
&\qquad =
4^{m+1}(d+\Qloss_m)^{m+3}(S+\Qloss_m)^{m+1}.
\end{aligned}
\]
Since \(\Qloss_m\le Q\) and \(d,S,Q\ge0\), monotonicity of powers gives the
claimed bound.
\end{proof}

\begin{remark}
The proof explains why a one-parameter total-defect envelope is the wrong
object for the balanced partial-transpose expansion.  The two defects pay two
different bases.  Collapsing them too early destroys the binomial structure
that produces \((d+Q)^{m+3}(S+Q)^{m+1}\).
\end{remark}

\section{One-Sided Genus Reduction}

The remaining task is to prove Assumption~\ref{ass:count}.  The reduction
used here passes through the one-sided defect \(\dplus\).

\begin{definition}[Plus-defect and genus counts]
Let
\[
  P_m(a)=\#\{\pi\in\Sym_{m+1}:\dplus(\pi)=a\}.
\]
For \(g\ge0\), define
\[
  G_m(g)=\#\{\pi\in\Sym_{m+1}:\dplus(\pi)=2g\}.
\]
\end{definition}

\begin{lemma}[Parity]\label{lem:parity}
For every \(\pi\in\Sym_{m+1}\), the plus-defect \(\dplus(\pi)\) is even.
\end{lemma}

\begin{proof}
Let \(n=m+1\).  The sign identity
\(\operatorname{sgn}(\gammacyc\pi)=\operatorname{sgn}(\gammacyc)
\operatorname{sgn}(\pi)\), together with
\(\operatorname{sgn}(\sigma)=(-1)^{n-\cyc(\sigma)}\), gives
\[
  n-\cyc(\gammacyc\pi)
  \equiv
  (n-1)+(n-\cyc(\pi))
  \pmod 2.
\]
Rearranging,
\[
  n+1-\cyc(\pi)-\cyc(\gammacyc\pi)\equiv0\pmod 2.
\]
This is exactly \(\dplus(\pi)\equiv0\pmod2\).
\end{proof}

\begin{assumption}[Genus-zero and slicing inputs]\label{ass:genus}
For every \(m\ge0\):
\begin{enumerate}[label=(\roman*),leftmargin=2.3em]
  \item the genus-zero count satisfies
  \[
    G_m(0)\le 4^{m+1};
  \]
  \item the genus slices satisfy the Chapuy-type recurrence
  \[
    G_m(g+1)\le (2(m+1))^3G_m(g)
    \qquad(g\ge0).
  \]
\end{enumerate}
\end{assumption}

\begin{proposition}[Iterated genus bound]\label{prop:genus}
Under Assumption~\ref{ass:genus},
\[
  G_m(g)\le 4^{m+1}(2(m+1))^{3g}
  \qquad(m,g\ge0).
\]
\end{proposition}

\begin{proof}
Induct on \(g\).  The case \(g=0\) is Assumption~\ref{ass:genus}(i).  If the
bound holds at \(g\), then Assumption~\ref{ass:genus}(ii) gives
\[
  G_m(g+1)\le (2(m+1))^3G_m(g)
  \le
  4^{m+1}(2(m+1))^{3(g+1)}.
\]
\end{proof}

\begin{proposition}[One-sided plus-count envelope]\label{prop:plus}
Under Assumption~\ref{ass:genus}, for every \(m,a\ge0\),
\[
  P_m(a)\le 4^{m+1}\Qloss_m^a.
\]
\end{proposition}

\begin{proof}
If \(a\) is odd, Lemma~\ref{lem:parity} gives \(P_m(a)=0\).  If
\(a=2g\), then \(P_m(a)=G_m(g)\), and Proposition~\ref{prop:genus} gives
\[
  P_m(a)\le 4^{m+1}(2(m+1))^{3g}.
\]
The chosen loss satisfies
\[
  \Qloss_m^2=(4(m+1))^{72}\ge (2(m+1))^3.
\]
Therefore
\[
  (2(m+1))^{3g}
  \le
  \Qloss_m^{2g}
  =
  \Qloss_m^a.
\]
The displayed inequality uses far less than the exponent \(36\); the larger
power is kept only to match the ambient Aubrun loss budget.
\end{proof}

\begin{theorem}[Genus route to the bidefect count]\label{thm:bidefect}
Assumption~\ref{ass:genus} implies Assumption~\ref{ass:count}.
\end{theorem}

\begin{proof}
Only the analytic ranges of Assumption~\ref{ass:count} are being claimed.
For fixed \(a,b\) in those ranges, the bidefect slice embeds in the
plus-defect slice:
\[
  \{\pi:\dplus(\pi)=a,\ \dminus(\pi)=b\}
  \subseteq
  \{\pi:\dplus(\pi)=a\}.
\]
Therefore
\[
  \Count_m(a,b)\le P_m(a)
  \le 4^{m+1}\Qloss_m^a.
\]
On the analytic ranges \(0\le a\le m+3\), \(0\le b\le m+1\), the binomial
coefficients are at least one, and \(\Qloss_m^a\le \Qloss_m^{a+b}\).
Thus
\[
  \Count_m(a,b)
  \le
  4^{m+1}\binom{m+3}{a}\binom{m+1}{b}\Qloss_m^{a+b}
  =
  \Envelope_m(a,b).
\]
\end{proof}

\begin{theorem}[Main bidefect reduction]\label{thm:main}
Under Assumption~\ref{ass:genus}, for every \(m\ge0\), every \(d,S,Q\ge0\)
with \(\Qloss_m\le Q\), the bidefect sum truncated to the analytic ranges
\(0\le a\le m+3\), \(0\le b\le m+1\) satisfies
\[
\sum_{a=0}^{m+3}\sum_{b=0}^{m+1}
  \Count_m(a,b)d^{m+3-a}S^{m+1-b}
  \le
  4^{m+1}(d+Q)^{m+3}(S+Q)^{m+1}.
\]
\end{theorem}

\begin{proof}
Combine Theorems~\ref{thm:bidefect} and~\ref{thm:scalar}.
\end{proof}

\begin{remark}
The theorem is only a statement over the analytic ranges displayed in the
double sum.  It makes no assertion about defect pairs with \(a>m+3\) or
\(b>m+1\).
\end{remark}

\section{Route to the Genus-Zero and Chapuy-Slicing Inputs}

This section records the enumerative inputs still needed to discharge
Assumption~\ref{ass:genus}.  It does not prove those inputs.  Its purpose is
to identify exactly where the remaining combinatorial work lies.

\subsection{Genus zero}

The condition \(\dplus(\pi)=0\) is the equality case of the Cayley inequality
\[
  \cyc(\pi)+\cyc(\gammacyc\pi)\le n+1.
\]
Equivalently, \(\pi\) lies on a geodesic between the identity and the long
cycle, up to the left/right convention.  The required genus-zero input is to
identify this equality case with noncrossing partitions of the linearly
ordered set \(\{0,\ldots,n-1\}\).  The desired bound
\[
  G_m(0)\le 4^{m+1}
\]
is then a Catalan-type estimate: noncrossing partitions of an \(n\)-point set
are counted by the Catalan number \(C_n\), and \(C_n\le4^n\).

One possible route is to prove noncrossing of the cycle partition and an
injectivity statement saying that the ordered predecessor data reconstructs
the original permutation.  Once that encoding is available, the separate
cardinality estimate for noncrossing partitions is elementary.

\subsection{Positive genus}

For \(\dplus(\pi)=2g\), the permutation pair \((\pi,\gammacyc\pi)\) should be
identified with the standard encoding of a unicellular map of genus \(g\).  The
missing enumerative input is a slicing map of Chapuy type: from genus \(g+1\)
one selects boundedly many marked local data and cuts to genus \(g\).  The
desired coarse form is only
\[
  G_m(g+1)\le (2n)^3G_m(g),
  \qquad n=m+1.
\]
This is far weaker than the sharp identities available for unicellular maps
\cite{Chapuy2011}, but it must be proved in the exact permutation model used
by the moment expansion.

\section{Relation to the Partial-Transpose Moment Expansion}

Aubrun's partially transposed Wishart moment formula assigns to each
\(\pi\in\Sym_{m+1}\) a contribution whose powers of the two large bases are
governed by the cycle counts appearing above.  The rank identity
Proposition~\ref{prop:rank} supplies the total exponent budget
\[
  \rank(\pi)=2m+4-\dplus(\pi)-\dminus(\pi).
\]
If the two bases are kept separate, a term of bidefect \((a,b)\) has the
schematic size
\[
  d^{m+3-a}S^{m+1-b}
\]
up to the universal polynomial loss \(\Qloss_m^{a+b}\) already present in
Aubrun's compatible-couple bound.  The two exponents are chosen so that their
sum recovers exactly the rank identity:
\[
  (m+3-a)+(m+1-b)=2m+4-a-b=\rank(\pi)
\]
when \(\dplus(\pi)=a\) and \(\dminus(\pi)=b\).  Thus the separated
bookkeeping keeps the two analytic bases visible while preserving the total
cycle-rank budget.  Theorem~\ref{thm:main} then gives the exact scalar
estimate needed by the shifted-base route.

This note does not reprove the whole random-matrix theorem.  Its contribution
is the clean combinatorial reduction: once the genus-zero count and Chapuy
slicing recurrence are supplied in the exact permutation model, the bidefect
count and scalar absorption follow with no further probabilistic input.

\section*{Verification Checklist for the Reduction}
\addcontentsline{toc}{section}{Verification Checklist for the Reduction}

A referee can verify the note by checking the following items.

\begin{enumerate}[leftmargin=2.3em]
  \item The oriented defects are nonnegative consequences of the two Cayley
  triangle inequalities.
  \item The rank identity is a direct algebraic identity.
  \item The scalar bound is exactly the binomial theorem applied twice.
  \item The plus-defect is even by the sign identity.
  \item The genus induction proves the displayed one-sided count envelope.
  \item The bidefect count follows from the one-sided count by inclusion and
  by the positivity of the two binomial coefficients on the analytic ranges.
  \item The only genuine remaining inputs are the genus-zero count and the
  Chapuy slicing recurrence in the exact permutation model.
\end{enumerate}

\begin{thebibliography}{9}

\bibitem{Aubrun2012}
G.~Aubrun.
\newblock Partial transposition of random states and non-centered
semicircular distributions.
\newblock \emph{Random Matrices: Theory and Applications} 1, 1250001,
2012.

\bibitem{AubrunSzarekYe2014}
G.~Aubrun, S.~J. Szarek, and D.~Ye.
\newblock Entanglement thresholds for random induced states.
\newblock \emph{Communications on Pure and Applied Mathematics} 67,
129--171, 2014.

\bibitem{Biane1997}
P.~Biane.
\newblock Some properties of crossings and partitions.
\newblock \emph{Discrete Mathematics} 175, 41--53, 1997.

\bibitem{Chapuy2011}
G.~Chapuy.
\newblock A new combinatorial identity for unicellular maps, via a direct
bijective approach.
\newblock \emph{Advances in Applied Mathematics} 47, 874--893, 2011.

\bibitem{CollinsNechita2010}
B.~Collins and I.~Nechita.
\newblock Random quantum channels I: graphical calculus and the Bell state
phenomenon.
\newblock \emph{Communications in Mathematical Physics} 297, 345--370,
2010.

\bibitem{FukudaSniady2013}
M.~Fukuda and P.~Sniady.
\newblock Partial transpose of random quantum states: exact formulas and
meanders.
\newblock \emph{Journal of Mathematical Physics} 54, 042202, 2013.

\end{thebibliography}

\end{document}
